%%%% ssac27-abstract.tex
%%%% Standalone abstract for the SSAC27 Research Paper Competition.
%%%%
%%%% SUBMISSION SPEC (verified at sloansportsconference.com, 2026-08-19):
%%%%   - Due Oct 1, 2026 11:59pm EST via Google Form.
%%%%   - Under 500 words INCLUDING title and body. Tables sit outside that count.
%%%%   - At most two tables and figures combined. We use one table.
%%%%   - Required sections: Introduction, Methods, Results, Conclusion.
%%%%   - Judged on novelty, academic rigor, reproducibility, and application.
%%%%   - Accepted manuscripts carry NO abstract section, so this document is a
%%%%     standalone selection artifact and must read without the paper.
%%%%   - Finalist title blocks show track and Paper ID with no author names; an
%%%%     ID is assigned after acceptance, so only the track appears here.
%%%%
%%%% THIS VERSION (2026-09-19): Prof. Highley's draft of 2026-09-15, typeset as written,
%%%% with three additions and nothing else changed:
%%%%   1. the four required headings;
%%%%   2. the numbers in the last paragraph, from the 25-season run
%%%%      (30 leagues x 25 seasons x 9 mechanisms = 6,750 seasons);
%%%%   3. the table (outside the count); the column and mechanism notes went to
%%%%      the delivery email, for Prof. Highley to add back if he wants them.
%%%% The last paragraph's counts use two measurement choices, both open to
%%%% revision: parity scored by the number of teams that never win a playoff
%%%% series in the run, and robustness scored by the exact counterfactual
%%%% (fifth-worst team dropped to worst, everything else held). Under those,
%%%% Simple and Beckett are undominated; Simple's help gains clear their paired
%%%% intervals against both baselines, its parity differences do not (except the
%%%% longest drought against 3-2-1), and the text says so. Every number
%%%% regenerates from the public repo:
%%%%   node testbed/harness/frontier_analysis.js --parity=neverWon --robust=causal
%%%% Prof. Highley's closing sentence about the frontier was moved under Conclusion so the
%%%% required heading is not empty; the one or two closing sentences he planned
%%%% go there.
%%%%
%%%% Word count: python3 wordcount_abstract.py

\documentclass[10pt]{article}

\usepackage[utf8]{inputenc}
\usepackage[T1]{fontenc}
\usepackage{times}
\usepackage[margin=0.6in,top=0.4in,bottom=0.4in]{geometry}
\usepackage{booktabs}
\usepackage{array}
\usepackage[table]{xcolor}
\usepackage[font=small,skip=3pt,justification=raggedright]{caption}
\usepackage{titlesec}
\usepackage{enumitem}
\usepackage[hidelinks]{hyperref}
\usepackage{url}
\urlstyle{same}

\definecolor{accent}{RGB}{234,240,247}

\titleformat{\section}{\normalfont\bfseries}{}{0pt}{}
\titlespacing*{\section}{0pt}{5pt}{2pt}

\setlength{\parindent}{0pt}
\setlength{\parskip}{2pt}
\setlist[itemize]{nosep, leftmargin=1.2em, topsep=1pt}

\pagestyle{empty}

\begin{document}

\begin{center}
{\LARGE\bfseries Carry-Over Lottery Allocation:\\[2pt] Practical Incentive-Compatible Drafts}\\[7pt]
{\normalsize Paper Track: Basketball}\\[2pt]
{\small Open-source repository: \url{https://github.com/kvr06-ai/cola-sloan-ssac27}}
\end{center}

\section{Introduction}

How can a league help weak teams through the draft without encouraging teams to tank? Carry-Over Lottery Allocation (COLA), presented here, is a family of draft mechanisms that all simultaneously discourage tanking and help the weakest teams.

Under any mechanism in the COLA framework, each team has a draft priority that rises and falls over time, but it does not rise based on more losses. That eliminates the tanking incentive. Draft priority carries over from year to year, being reduced when a team has playoff success or makes a high draft pick. That directs the most help to the weakest teams over time.

For 40 years, the NBA has navigated the tradeoff between helping weak teams and discouraging tanking. COLA solves that primary dilemma, but there are other tradeoffs worth considering among the COLA variants. After demonstrating COLA's incentive compatibility, this paper discusses the tradeoffs between mechanisms within the framework.

\section{Methods}

We identify several tunable dials for mechanisms within the COLA framework:
\begin{itemize}
\item Which teams are eligible for the lottery?
\item How does a team's lottery priority increase?
\item How does the priority decrease?
\item How does the priority determine lottery odds and draft positions?
\end{itemize}

We consider different COLA variants plus other draft mechanisms, including 3-2-1 Lottery, and compare them using ZenGM Basketball-GM, an open-source simulator that plays out every game and roster move. We measure the mechanisms against these criteria:
\begin{itemize}
\item \textbf{Targeted help.} Does help reach the weakest teams, as measured by the difference between the help received by the highest priority team and other teams?
\item \textbf{Long-term parity.} Can all teams be competitive, as measured by the longest drought any team has between playoff series wins?
\item \textbf{Simplicity.} How complicated is the mechanism? Though measured subjectively, it is not difficult to identify complexity.
\item \textbf{Robustness.} Does the mechanism admit tanking incentives at the margin? COLA variants are incentive-compatible under reasonable assumptions. Further relaxation of those assumptions may entail a choice between extra complexity or allowing potential tanking incentives in rare cases.
\end{itemize}

\section{Results}

Many mechanisms are dominated under these criteria, meaning there is another mechanism that is better on all four.

We ran every mechanism on the same 30 simulated leagues for 25 seasons each, 6,750 seasons in all. Because the leagues are shared, we compare mechanisms by paired differences with 95\% confidence intervals. Simple COLA helps long-drought teams more than the current NBA lottery does, and more than 3-2-1 does. Both gains clear their confidence intervals. On parity, Simple COLA does at least as well as both. It never reads a team's record, so it offers no reward for losing. Simple COLA or Beckett COLA matches or beats every other mechanism in every column of Table~1. Parity is where every mechanism fell short. In every league some team went close to 25 seasons without winning a playoff series. No mechanism cut the number of never-winners by more than chance could explain.
% Filled from the 25-season run. "2019--2026": the June 2026 draft also ran under
% the old rules; 3-2-1 first applies to the 2027 draft. Sources for each claim
% (runs/frontier/analysis_recut.txt, --parity=neverWon --robust=causal):
%   help vs 2019 lottery +0.17 [0.13, 0.21] places per drought season, vs 3-2-1
%   +0.22 [0.19, 0.26]; longest drought vs 2019 lottery -0.07 [-1.06, 0.93] and
%   never-winners -0.03 [-0.56, 0.49]; longest drought vs 3-2-1 -1.10 [-2.01,
%   -0.19] and never-winners -0.60 [-1.21, 0.01]; reward for losing 0 by
%   construction (tank_counterfactual.js gives the 2019 lottery 1.32).

\begin{center}
\small
\setlength{\tabcolsep}{7pt}
\renewcommand{\arraystretch}{1.0}
\begin{tabular}{l r r r r}
\toprule
\textbf{Mechanism} & \textbf{Help to} & \textbf{Reward for} & \textbf{Teams never} & \textbf{Longest} \\
 & \textbf{long-drought teams} & \textbf{losing} & \textbf{winning a series} & \textbf{drought} \\
 & \textit{draft places per} & \textit{draft places gained} & \textit{in 25 years,} & \textit{seasons,} \\
 & \textit{drought season} & \textit{by finishing last} & \textit{per league} & \textit{of 25} \\
\midrule
Current NBA lottery (2019 to 2026) & $0.07$ & $1.32$ & $0.63$ & $23.3$ \\
New NBA lottery, 3-2-1 (from 2027) & $0.02$ & losing costs picks & $1.20$ & $24.4$ \\
Equal-odds lottery & $0.00$ & $0$ & $0.93$ & $23.7$ \\
Classic COLA, the original version & $0.10$ & $1.85$ & $0.93$ & $24.1$ \\
Waitlist COLA & $0.13$ & $0$ & $1.10$ & $24.4$ \\
\rowcolor{accent}
Simple COLA & $0.24$ & $0$ & $0.60$ & $23.3$ \\
Full-lottery COLA & $0.09$ & $0$ & $0.83$ & $23.8$ \\
Countdown COLA & $0.19$ & losing costs picks & $1.33$ & $24.6$ \\
Beckett COLA & $0.28$ & losing costs picks & $0.83$ & $23.7$ \\
\bottomrule
\end{tabular}\\[2pt]
\begin{minipage}{0.99\linewidth}\footnotesize
\textbf{Table 1.} Nine mechanisms on the same 30 simulated leagues of 25 seasons each.
\end{minipage}
\end{center}

\section{Conclusion}

We identify a frontier of mechanisms that are not dominated and therefore viable options for implementation, depending on the league's goals.
% One or two closing sentences to add here (Prof. Highley, 2026-09-15).

\end{document}
